\thispagestyle{empty}
\centering {\Huge \textbf{ABSTRACT}}
\vskip 2.5mm
\rule{\textwidth}{0.8pt}\\[-3    mm]
\rule{\textwidth}{0.8pt}
\vskip 2cm
\justifying
\begin{tcolorbox}[colback=white, colframe=black, boxrule=0.6pt, arc=0pt, left=8pt, right=8pt, top=6pt, bottom=6pt]
    \hspace{6mm} This Master's thesis investigates the dynamics of \textbf{open quantum systems subjected to continuous measurement of their environment}, within the framework of quantum trajectories. Starting from the \textbf{input-output formalism}, we derive the Lindblad master equation under the Born-Markov approximation, describing the unconditional (unmonitored) evolution of a system coupled to a continuum of bosonic environmental modes. We then consider the case where the environment is continuously monitored, focusing on two fundamental detection schemes: the \textbf{photo-detection} and the \textbf{homodyne detection}. For each scheme we derive the corresponding stochastic master equation governing the \textbf{conditional state} of the system. We show that photo-detection yields a jump-like unravelling driven by a Poisson process, while homodyne detection yields a diffusive unravelling driven by a Wiener process. We further discuss the linearization of these SMEs and their reformulation in terms of Kraus operators, which provides the bases for numerical simulation. These theoretical results are illustrated by numerical simulations, implemented in \verb|Julia| using the \verb|QuantumToolbox.jl| package. In particular, we simulate the dynamics of a qubit and of a harmonic oscillator coupled to a zero-temperature bath. For both detection schemes, individual quantum trajectories are simulated and averaged over a large number of realizations. The results confirm that \textbf{the average over trajectories reproduces the unconditional Lindblad dynamics}, consistent with the unravelling interpretation.
\end{tcolorbox}
\vspace{1cm}\centering\rule{0.5\textwidth}{1pt}\vspace{1cm}
\begin{tcolorbox}[colback=white, colframe=black, boxrule=0.6pt, arc=0pt, left=8pt, right=8pt, top=6pt, bottom=6pt]
    \hspace{6mm} Ce mémoire de Master examine la dynamique des \textbf{systèmes quantiques ouverts soumis à une mesure continue de leur environnement}, inscrit dans le cadre des trajectoires quantiques. En partant du \textbf{formalisme entrée-sortie} (input-output), nous dérivons l'équation maîtresse de Lindblad sous l'approximation de Born-Markov, qui décrit l'évolution inconditionnelle (sans mesure continue) d'un système couplé à un continuum de modes environnementaux bosoniques. Nous considérons ensuite le cas où l'environnement est mesuré en continu, en nous concentrant sur deux schémas de détection fondamentaux : la \textbf{photodétection} et la \textbf{détection homodyne}. Pour chaque cas de mesure, nous dérivons l'équation stochastique maîtresse correspondante régissant l'\textbf{état conditionnel} du système. Nous montrons que la photodétection conduit à un ``dénouement'' (unravelling) par sauts, piloté par un processus de Poisson, tandis que la détection homodyne donne lieu à un dénouement diffusif, dirigé par un processus de Wiener. Nous présentons de surcroît la linéarisation de ces équations stochastiques maîtresses et leur reformulation en termes d'opérateurs de Kraus, ce qui fournit les bases de la simulation numérique. Ces résultats théoriques sont illustrés par des simulations numériques implémentées en \verb|Julia| à l'aide du package \verb|QuantumToolbox.jl|. En particulier, nous simulons la dynamique d'un qubit et d'un oscillateur harmonique couplés à un bain de température nulle. Pour les deux schémas de détection, les trajectoires quantiques individuelles sont simulées et moyennées sur un grand nombre de réalisations. Les résultats confirment que \textbf{la moyenne des trajectoires reproduit la dynamique inconditionnelle de Lindblad}, ce qui est cohérent avec l'interprétation par du dénouement des trajectoires.
\end{tcolorbox}